\documentclass[12pt, letterpaper]{article}

\usepackage[utf8]{inputenc}
\usepackage[T1]{fontenc}
\usepackage[spanish]{babel}
\usepackage{geometry}
\usepackage{amsmath, amssymb, amsthm}
\usepackage{booktabs}
\usepackage{array}
\usepackage{microtype}
\usepackage{setspace}
\usepackage{parskip}
\usepackage{titlesec}
\usepackage{fancyhdr}
\usepackage{enumitem}
\usepackage{bm}
\usepackage{mathtools}
\usepackage{tcolorbox}
\usepackage{listings}
\usepackage{xcolor}
\usepackage{hyperref}

\tcbuselibrary{skins, breakable}

\geometry{top=2.8cm, bottom=3.0cm, left=3.2cm, right=3.2cm, headheight=14pt}
\setstretch{1.20}

% ---- Encabezado y pie de página ----
\pagestyle{fancy}
\fancyhf{}
\fancyhead[C]{\small\textsc{Econometría --- Gujarati --- Ejercicio 13.4}}
\fancyfoot[C]{\small\thepage}
\renewcommand{\headrulewidth}{0.4pt}
\renewcommand{\footrulewidth}{0pt}

% ---- Títulos de sección ----
\titleformat{\section}
  {\large\bfseries\scshape}{\thesection.}{0.6em}{}
  [\vspace{-4pt}\rule{\linewidth}{0.4pt}]
\titleformat{\subsection}{\normalsize\bfseries}{\thesubsection.}{0.5em}{}
\titleformat{\subsubsection}{\normalsize\itshape}{}{0em}{}

% ---- Caja resaltada ----
\newtcolorbox{clave}{
  colback=white, colframe=black,
  boxrule=0.5pt, arc=3pt,
  left=8pt, right=8pt, top=6pt, bottom=6pt,
  breakable
}

% ---- Estilo de código Stata ----
\definecolor{stataback}{RGB}{248,248,248}
\definecolor{stataframe}{RGB}{180,180,180}
\definecolor{statacomment}{RGB}{100,100,100}
\definecolor{statakeyword}{RGB}{0,70,140}

\lstdefinelanguage{Stata}{
  morekeywords={clear, set, gen, generate, quietly, regress, scalar,
                summarize, di, display, estimates, store, table,
                correlate, test, estat, log, close, seed, obs,
                noconstant, as, text, result, newline, if, replace},
  sensitive=false,
  morecomment=[l]{*},
  morecomment=[l]{//},
  morestring=[b]",
}

\lstset{
  language=Stata,
  backgroundcolor=\color{stataback},
  basicstyle=\ttfamily\footnotesize,
  keywordstyle=\color{statakeyword}\bfseries,
  commentstyle=\color{statacomment}\itshape,
  stringstyle=\color{black},
  frame=single,
  framesep=4pt,
  rulecolor=\color{stataframe},
  xleftmargin=8pt,
  xrightmargin=4pt,
  breaklines=true,
  breakatwhitespace=true,
  showstringspaces=false,
  tabsize=4,
  captionpos=b,
  numbers=none,
}

\newcommand{\Var}{\operatorname{Var}}

% ---- Entorno para salida de Stata ----
\lstdefinestyle{stataout}{
  language={},
  backgroundcolor=\color{white},
  basicstyle=\ttfamily\footnotesize,
  frame=single,
  framesep=4pt,
  rulecolor=\color{stataframe},
  xleftmargin=8pt,
  xrightmargin=4pt,
  breaklines=true,
  showstringspaces=false,
}

% ===========================================================================
\begin{document}
% ===========================================================================

\begin{center}
  {\LARGE\bfseries Ejercicio 13.4 --- Inclusión de una Variable Irrelevante}\\[6pt]
  {\LARGE\bfseries en el Modelo de Regresión Múltiple}\\[12pt]
  {\large\itshape Consecuencias del error de especificación por inclusión}\\
  {\large\itshape de una variable cuyo verdadero coeficiente es cero}\\[14pt]
  \rule{0.55\linewidth}{0.6pt}\\[8pt]
  {\normalsize Ejercicio~13.4 --- \textit{Econometría}, Gujarati \& Porter, 5.a ed.\
  por Emanuel Quintana Silva}\\[4pt]
  {\small\textsc{Verificación empírica con datos simulados, $n = 200$, semilla 2026}}
\end{center}

\vspace{10pt}

%------------------------------------------------------------------------------
\section{Planteamiento del problema}
%------------------------------------------------------------------------------

Suponga que el verdadero modelo es:

\begin{align}
  \textbf{Modelo verdadero:} \quad
  Y_i &= \beta_1 + \beta_2 X_{2i} + u_i \label{eq:verdadero}\\[6pt]
  \textbf{Modelo mal especificado:} \quad
  Y_i &= \alpha_1 + \alpha_2 X_{2i} + \alpha_3 X_{3i} + v_i \label{eq:ampliado}
\end{align}

donde $X_3$ es una variable \textbf{irrelevante} en el sentido de que su
verdadero coeficiente poblacional es $\beta_3 = 0$. El modelo~\eqref{eq:ampliado}
incluye esta variable sin que la teoría económica la justifique.

En la taxonomía de Gujarati (sección~13.2), este error corresponde a la
\textbf{inclusión de una variable irrelevante} o \textit{sobreajuste del modelo}.
La comparación con el ejercicio~13.3 es conceptualmente decisiva:

\begin{center}
\renewcommand{\arraystretch}{1.3}
\small
\begin{tabular}{lll}
\toprule
\textbf{Tipo de error} & \textbf{Consecuencia principal} & \textbf{Ejercicio} \\
\midrule
Omisión de variable relevante & Sesgo + inconsistencia & 13.3 \\
\textbf{Inclusión de variable irrelevante} & \textbf{Insesgadez + ineficiencia} & \textbf{13.4} \\
\bottomrule
\end{tabular}
\end{center}

%------------------------------------------------------------------------------
\section{Consecuencias teóricas}
%------------------------------------------------------------------------------

\subsection{Inciso a) --- Comportamiento de $R^2$ y $\bar{R}^2$}

\subsubsection{$R^2$ convencional (no penalizado)}

El $R^2$ del modelo~\eqref{eq:ampliado} será \textbf{mayor o al menos igual}
al del modelo~\eqref{eq:verdadero}. Este resultado es un \textbf{hecho
algebraico}, no estadístico:

\begin{equation}
  R^2 = 1 - \frac{SCR}{SCT} = 1 - \frac{\sum \hat{v}_i^2}{\sum(Y_i - \bar{Y})^2}
\end{equation}

La Suma de Cuadrados de los Residuales ($SCR = \sum \hat{v}_i^2$) \textbf{nunca
puede aumentar} al añadir un regresor adicional a un modelo MCO, porque el
conjunto de parámetros del modelo~\eqref{eq:verdadero} es un subconjunto de
los de~\eqref{eq:ampliado}. En el peor caso, el estimador del nuevo
coeficiente puede ser exactamente cero, replicando la $SCR$ original.
Formalmente:

\begin{equation}
  SCR_{\eqref{eq:ampliado}} \leq SCR_{\eqref{eq:verdadero}}
  \implies
  \boxed{R^2_{\eqref{eq:ampliado}} \geq R^2_{\eqref{eq:verdadero}}}
\end{equation}

Esta propiedad es la razón por la que el $R^2$ convencional constituye una
medida de bondad de ajuste \textbf{deficiente} para comparar modelos con
distinto número de regresores: siempre premia la inclusión de variables,
independientemente de su relevancia teórica o estadística.

\subsubsection{$R^2$ ajustado o corregido ($\bar{R}^2$)}

El $\bar{R}^2$ del modelo~\eqref{eq:ampliado} \textbf{probablemente será menor
o igual} al del modelo~\eqref{eq:verdadero}. A diferencia del $R^2$
convencional, el ajustado incorpora una penalización explícita por el número
de parámetros estimados:

\begin{equation}
  \bar{R}^2 = 1 - \frac{SCR/(n-k)}{SCT/(n-1)}
\end{equation}

donde $k$ es el número de parámetros (incluyendo el intercepto). Al añadir
$X_3$, el denominador de $SCR/(n-k)$ disminuye de $n-2$ a $n-3$, lo que
tiende a aumentar la varianza estimada del error aunque $SCR$ se reduzca
marginalmente. Puede demostrarse que el $\bar{R}^2$ aumenta si, y solo si,
el valor absoluto del estadístico $t$ de la variable añadida supera la
unidad:

\begin{equation}
  \bar{R}^2_{\eqref{eq:ampliado}} > \bar{R}^2_{\eqref{eq:verdadero}}
  \iff |t_{\hat{\alpha}_3}| > 1
\end{equation}

Dado que el verdadero $\beta_3 = 0$, el estadístico $t$ de $\hat{\alpha}_3$
tenderá a ser pequeño. Por lo tanto:

\begin{equation}
  \boxed{\bar{R}^2_{\eqref{eq:ampliado}} \leq \bar{R}^2_{\eqref{eq:verdadero}}}
\end{equation}

\begin{center}
\renewcommand{\arraystretch}{1.3}
\small
\begin{tabular}{p{2.8cm}p{3.2cm}p{3.8cm}p{3.4cm}}
\toprule
\textbf{Medida} & \textbf{Modelo \eqref{eq:verdadero}} & \textbf{Modelo \eqref{eq:ampliado}} & \textbf{Relación esperada} \\
\midrule
$R^2$       & $R^2_{(1)}$       & $R^2_{(2)}$       & $R^2_{(2)} \geq R^2_{(1)}$       \\
$\bar{R}^2$ & $\bar{R}^2_{(1)}$ & $\bar{R}^2_{(2)}$ & $\bar{R}^2_{(2)} \leq \bar{R}^2_{(1)}$ \\
\bottomrule
\end{tabular}
\end{center}

\subsection{Inciso b) --- Insesgadez y consistencia de los estimadores}

\textbf{Sí: los estimadores MCO del modelo~\eqref{eq:ampliado} son
insesgados y consistentes.} Esta es la propiedad central que distingue la
sobre-especificación de la sub-especificación.

\subsubsection{Demostración algebraica}

El estimador MCO del vector de coeficientes en el modelo~\eqref{eq:ampliado} es:

\begin{equation}
  \hat{\boldsymbol{\alpha}} = (\mathbf{X}'\mathbf{X})^{-1}\mathbf{X}'\mathbf{Y}
\end{equation}

donde $\mathbf{X} = [\mathbf{1},\; \mathbf{X}_2,\; \mathbf{X}_3]$.
Sustituyendo el modelo verdadero
$\mathbf{Y} = \beta_1\mathbf{1} + \beta_2\mathbf{X}_2 + \mathbf{u}$:

\begin{equation}
  \hat{\boldsymbol{\alpha}}
  = \beta_1(\mathbf{X}'\mathbf{X})^{-1}\mathbf{X}'\mathbf{1}
  + \beta_2(\mathbf{X}'\mathbf{X})^{-1}\mathbf{X}'\mathbf{X}_2
  + (\mathbf{X}'\mathbf{X})^{-1}\mathbf{X}'\mathbf{u}
\end{equation}

Tomando la esperanza y usando el supuesto clásico
$E(\mathbf{u}\mid\mathbf{X}) = \mathbf{0}$:

\begin{equation}
  E(\hat{\boldsymbol{\alpha}})
  = \beta_1(\mathbf{X}'\mathbf{X})^{-1}\mathbf{X}'\mathbf{1}
  + \beta_2(\mathbf{X}'\mathbf{X})^{-1}\mathbf{X}'\mathbf{X}_2
\end{equation}

Nótese que $\mathbf{X}'\mathbf{1}$ y $\mathbf{X}'\mathbf{X}_2$ son,
respectivamente, la primera y la segunda columna de $\mathbf{X}'\mathbf{X}$,
de modo que:

\begin{equation}
  (\mathbf{X}'\mathbf{X})^{-1}\mathbf{X}'\mathbf{1} = \mathbf{e}_1,
  \qquad
  (\mathbf{X}'\mathbf{X})^{-1}\mathbf{X}'\mathbf{X}_2 = \mathbf{e}_2
\end{equation}

donde $\mathbf{e}_j$ denota el $j$-ésimo vector canónico. Por tanto:

\begin{equation}
  E(\hat{\boldsymbol{\alpha}}) = \beta_1\mathbf{e}_1 + \beta_2\mathbf{e}_2
  =
  \begin{pmatrix}\beta_1 \\ \beta_2 \\ 0\end{pmatrix}
\end{equation}

Lo que confirma que:

\begin{equation}
  \boxed{
    E(\hat{\alpha}_1) = \beta_1, \qquad
    E(\hat{\alpha}_2) = \beta_2, \qquad
    E(\hat{\alpha}_3) = 0
  }
\end{equation}

Los tres estimadores son insesgados: los dos primeros recuperan los
verdaderos parámetros del modelo, y el tercero estima correctamente su
valor poblacional (que es cero).

\subsubsection{Intuición económica}

La insesgadez se preserva porque los supuestos del Modelo Clásico de
Regresión Lineal (MCRL) se mantienen intactos en el modelo~\eqref{eq:ampliado}.
La única violación sería si $X_3$ estuviese correlacionada con el término
de error $u_i$; pero como $X_3$ es exógena y su verdadero coeficiente es
simplemente cero (no que la variable ``no exista''), la condición de
exogeneidad no se viola.

\subsubsection{Consistencia}

Dado que los estimadores son insesgados y sus varianzas tienden a cero
cuando $n \to \infty$, también son \textbf{consistentes}:

\begin{equation}
  \hat{\alpha}_j \xrightarrow{p} \beta_j \qquad \text{cuando } n \to \infty
\end{equation}

\begin{clave}
  A diferencia del Ej.~13.3, donde $E(\hat{\beta}_1) \neq \alpha_1$ de
  forma sistemática, aquí la insesgadez \textbf{se preserva}. Este es el
  contraste fundamental entre sub-especificación (sesgo) y sobre-especificación
  (ineficiencia sin sesgo).
\end{clave}

\subsection{Inciso c) --- Efecto sobre las varianzas de $\hat{\beta}_1$ y $\hat{\beta}_2$}

\textbf{Sí: las varianzas en general aumentan. Esta es la consecuencia
estadística más importante del sobreajuste.}

\subsubsection{Demostración general}

La matriz de varianzas y covarianzas de los estimadores MCO es:

\begin{equation}
  \Var(\hat{\boldsymbol{\alpha}}) = \sigma^2 (\mathbf{X}'\mathbf{X})^{-1}
\end{equation}

Al añadir $X_3$, la matriz $\mathbf{X}'\mathbf{X}$ se expande de dimensión
$2 \times 2$ a $3 \times 3$. En general, los elementos diagonales de
$(\mathbf{X}'\mathbf{X})^{-1}$ no pueden disminuir cuando se añaden
columnas a $\mathbf{X}$, de modo que:

\begin{equation}
  \Var(\hat{\alpha}_j) \geq \Var(\hat{\beta}_j), \quad j = 1, 2
\end{equation}

con igualdad únicamente en el caso límite de ortogonalidad perfecta entre
$X_2$ y $X_3$.

\subsubsection{El Factor de Inflación de la Varianza (VIF)}

Para una regresión con dos regresores ($X_2$ y $X_3$), la varianza del
estimador de la pendiente en el modelo~\eqref{eq:ampliado} es:

\begin{equation}
  \Var(\hat{\alpha}_2)
  = \frac{\sigma^2}{\sum x_{2i}^2 \,(1 - r_{23}^2)}
  = \underbrace{\frac{\sigma^2}{\sum x_{2i}^2}}_{\Var(\hat{\beta}_2)}
    \cdot
    \underbrace{\frac{1}{1 - r_{23}^2}}_{\mathrm{VIF}}
\end{equation}

donde $r_{23}$ es el coeficiente de correlación muestral entre $X_2$ y
$X_3$, y $x_{2i} = X_{2i} - \bar{X}_2$. La razón de varianzas define el
Factor de Inflación de la Varianza:

\begin{equation}
  \boxed{
    \mathrm{VIF} = \frac{\Var(\hat{\alpha}_2)}{\Var(\hat{\beta}_2)}
    = \frac{1}{1 - r_{23}^2} \geq 1
  }
\end{equation}

El VIF mide cuántas veces se infla la varianza del estimador por efecto de
la correlación con el regresor adicional.

\subsubsection{Análisis por escenarios de correlación}

\begin{center}
\renewcommand{\arraystretch}{1.3}
\small
\begin{tabular}{cccp{5.8cm}}
\toprule
\textbf{Correlación $r_{23}$} & $r_{23}^2$ & \textbf{VIF} & \textbf{Efecto sobre $\Var(\hat{\beta}_2)$} \\
\midrule
0.000 & 0.000 &  1.00 & Sin cambio (caso límite de ortogonalidad) \\
0.300 & 0.090 &  1.10 & Aumento del 10\% \\
0.500 & 0.250 &  1.33 & Aumento del 33\% \\
0.700 & 0.490 &  1.96 & Aumento del 96\% \\
0.900 & 0.810 &  5.26 & Aumento del 426\% \\
0.977 & 0.955 & 22.01 & Aumento del 2101\% \quad (caso empírico) \\
0.990 & 0.980 & 50.25 & Aumento del 4925\% \\
\bottomrule
\end{tabular}
\end{center}

\textit{Caso especial}: si $r_{23} = 0$ (ortogonalidad perfecta), el
VIF $= 1$ y las varianzas no cambian. En datos económicos reales esta
condición prácticamente nunca se satisface.

\subsubsection{Consecuencias para la inferencia estadística}

Aunque los estimadores siguen siendo insesgados y las pruebas formalmente
válidas, la inflación de varianzas tiene efectos directos:

\begin{enumerate}[leftmargin=1.8em]
  \item \textbf{Errores estándar más grandes}: $SE(\hat{\alpha}_j) \geq
    SE(\hat{\beta}_j)$, lo que produce estadísticos $t$ más pequeños.
  \item \textbf{Intervalos de confianza más amplios}: el mismo nivel
    $(1-\alpha)$ requiere intervalos más anchos, reduciendo la precisión
    práctica.
  \item \textbf{Menor potencia de las pruebas}: es más difícil rechazar
    $H_0: \beta_j = 0$ aunque sea falsa (mayor probabilidad de Error
    Tipo~II).
  \item \textbf{Multicolinealidad inducida artificialmente}: si $X_3$
    está correlacionada con $X_2$, se introduce un problema de
    multicolinealidad que no existía en el modelo verdadero.
\end{enumerate}

%------------------------------------------------------------------------------
\section{Resumen de consecuencias}
%------------------------------------------------------------------------------

\begin{center}
\renewcommand{\arraystretch}{1.4}
\small
\begin{tabular}{p{3.8cm}p{3.8cm}p{3.8cm}}
\toprule
\textbf{Propiedad}
  & \textbf{Modelo \eqref{eq:verdadero}} \newline (correcto)
  & \textbf{Modelo \eqref{eq:ampliado}} \newline (sobre-especificado) \\
\midrule
Insesgadez             & Preservada                  & Preservada \\
Consistencia           & Preservada                  & Preservada \\
Eficiencia (MELI)      & Óptima (Gauss-Markov)       & Reducida: $\Var(\hat{\alpha}) \geq \Var(\hat{\beta})$ \\
Validez $t$ y $F$      & Válidas                     & Válidas (pero menos potentes) \\
$R^2$ convencional     & Valor base                  & Mayor o igual al de \eqref{eq:verdadero} \\
$\bar{R}^2$ ajustado   & Valor base                  & Menor o igual al de \eqref{eq:verdadero} \\
VIF                    & 1 (base)                    & $\geq 1$ (se infla según $r_{23}^2$) \\
Predicciones           & Insesgadas y eficientes     & Insesgadas pero menos precisas \\
\bottomrule
\end{tabular}
\end{center}

%------------------------------------------------------------------------------
\section{Verificación empírica con datos simulados}
%------------------------------------------------------------------------------

Se generan datos del proceso verdadero $Y_i = 2 + 3X_{2i} + u_i$
($\beta_1 = 2$, $\beta_2 = 3$, $\beta_3 = 0$, $u_i \sim N(0,1)$, $n = 200$)
y se añaden tres versiones de $X_3$ con distintos niveles de correlación
con $X_2$ para verificar el efecto VIF empíricamente.

\subsection{Código Stata --- simulación principal}

\begin{lstlisting}[caption={Simulación: modelo correcto vs. modelos sobre-especificados}]
clear
set seed 2026
set obs 200

gen X2 = runiform(1, 10)
gen u  = rnormal(0, 1)
gen Y  = 2 + 3 * X2 + u

* Variable irrelevante con correlacion baja  (~0.03)
gen X3_low = 0.05 * X2 + rnormal(0, 2)

* Variable irrelevante con correlacion moderada (~0.50)
gen X3_mid = 0.50 * X2 + rnormal(0, 1.5)

* Variable irrelevante con correlacion alta  (~0.977)
gen X3_hi  = 0.90 * X2 + rnormal(0, 0.5)

* ---- Modelo (1): correcto ----
quietly regress Y X2
scalar b2_correct = _b[X2]
scalar se_b2      = _se[X2]
scalar var_b2     = se_b2^2

* ---- Modelos (2a), (2b), (2c): sobre-especificados ----
quietly regress Y X2 X3_low
quietly regress Y X2 X3_mid
quietly regress Y X2 X3_hi

* ---- Tabla comparativa ----
quietly regress Y X2
estimates store m_correcto
quietly regress Y X2 X3_low
estimates store m_bajo
quietly regress Y X2 X3_mid
estimates store m_medio
quietly regress Y X2 X3_hi
estimates store m_alto

estimates table m_correcto m_bajo m_medio m_alto, ///
    b(%9.4f) se stats(N r2 r2_a)
\end{lstlisting}

\subsection{Resultados --- modelo correcto}

\begin{center}
\renewcommand{\arraystretch}{1.3}
\small
\begin{tabular}{lrrrr}
\toprule
\textbf{Variable} & \textbf{Coeficiente} & \textbf{Error est.} & \textbf{Razón $t$} & \textbf{$P(>|t|)$} \\
\midrule
$X_{2i}$  & 2.9797 & 0.0276 & 107.96 & 0.000 \\
Constante & 2.0436 & 0.1657 &  12.33 & 0.000 \\
\midrule
\multicolumn{5}{l}{$R^2 = 0.9833$\quad $\bar{R}^2 = 0.9832$\quad $n = 200$} \\
\bottomrule
\end{tabular}
\end{center}

\subsection{Resultados --- tabla comparativa de los cuatro modelos}

\begin{center}
\renewcommand{\arraystretch}{1.3}
\small
\begin{tabular}{lcccc}
\toprule
\textbf{Variable}
  & \textbf{(1) Correcto}
  & \textbf{(2a) $r_{23} \approx 0.03$}
  & \textbf{(2b) $r_{23} \approx 0.50$}
  & \textbf{(2c) $r_{23} \approx 0.977$} \\
\midrule
$X_2$ (coef.)       & 2.9797  & 2.9795  & 2.9858  & 2.8618  \\
$X_2$ (SE)          & 0.0276  & 0.0274  & 0.0392  & 0.1330  \\
\midrule
$X_{3,\text{low}}$  &         & $-0.0809$ &       &         \\
(SE)                &         & 0.0395    &       &         \\
$X_{3,\text{mid}}$  &         &           & $-0.0110$ &     \\
(SE)                &         &           & 0.0496    &     \\
$X_{3,\text{hi}}$   &         &           &       & 0.1296  \\
(SE)                &         &           &       & 0.1431  \\
\midrule
Constante           & 2.0436  & 2.0731  & 2.0399  & 2.0542  \\
(SE)                & 0.1657  & 0.1650  & 0.1670  & 0.1662  \\
\midrule
$n$                 & 200     & 200     & 200     & 200     \\
$R^2$               & 0.9833  & 0.9837  & 0.9833  & 0.9834  \\
$\bar{R}^2$         & 0.9832  & 0.9835  & 0.9831  & 0.9832  \\
VIF empírico        & 1.00    & 0.98    & 2.02    & 23.23   \\
\bottomrule
\end{tabular}
\end{center}

%------------------------------------------------------------------------------
\section{Verificación algebraica del VIF}
%------------------------------------------------------------------------------

\subsection{Código Stata}

\begin{lstlisting}[caption={Verificación: VIF teórico vs. VIF empírico}]
clear
set seed 2026
set obs 200

gen X2    = runiform(1, 10)
gen X3_hi = 0.90 * X2 + rnormal(0, 0.5)
gen u     = rnormal(0, 1)
gen Y     = 2 + 3 * X2 + u

* Correlacion muestral entre X2 y X3_hi
correlate X2 X3_hi
scalar r23     = r(rho)
scalar vif_teo = 1 / (1 - r23^2)

* Varianzas empiricas de la pendiente
quietly regress Y X2
scalar var_b2 = _se[X2]^2

quietly regress Y X2 X3_hi
scalar var_a2  = _se[X2]^2
scalar vif_emp = var_a2 / var_b2

di "Correlacion r(X2,X3):         " r23
di "r23^2:                        " r23^2
di "VIF teorico = 1/(1-r23^2):    " vif_teo
di "Var(beta_2)  modelo correcto: " var_b2
di "Var(alpha_2) modelo ampliado: " var_a2
di "VIF empirico (var_a2/var_b2): " vif_emp
di "Diferencia (emp - teo):       " vif_emp - vif_teo
\end{lstlisting}

\subsection{Resultados}

\begin{center}
\renewcommand{\arraystretch}{1.3}
\small
\begin{tabular}{lr}
\toprule
\textbf{Cantidad} & \textbf{Valor} \\
\midrule
Correlación $r_{23}$ entre $X_2$ y $X_{3,\text{hi}}$     & 0.9770   \\
$r_{23}^2$                                                & 0.9546   \\
VIF teórico $= 1/(1 - r_{23}^2)$                         & 22.0074  \\
\midrule
$\Var(\hat{\beta}_2)$ modelo correcto~\eqref{eq:verdadero}  & 0.000608 \\
$\Var(\hat{\alpha}_2)$ modelo ampliado~\eqref{eq:ampliado}  & 0.013157 \\
VIF empírico $= \Var(\hat{\alpha}_2)/\Var(\hat{\beta}_2)$ & 21.6570  \\
Diferencia (empírico $-$ teórico)                         & $-0.3504$ \\
\bottomrule
\end{tabular}
\end{center}

\begin{clave}
  El VIF empírico ($21.66$) replica con alta precisión el VIF teórico
  ($22.01$), con una diferencia de apenas $-0.35$. Esto confirma que la
  fórmula $\mathrm{VIF} = 1/(1 - r_{23}^2)$ describe correctamente la
  inflación de varianzas producida por la variable irrelevante.
\end{clave}

%------------------------------------------------------------------------------
\section{Diagnóstico empírico recomendado}
%------------------------------------------------------------------------------

En la práctica, para detectar si una variable incluida es irrelevante se
recurre a las siguientes herramientas:

\subsection{Código Stata}

\begin{lstlisting}[caption={Diagnóstico: VIF, prueba F parcial y criterios de información}]
clear
set seed 2026
set obs 200
gen X2     = runiform(1, 10)
gen u      = rnormal(0, 1)
gen Y      = 2 + 3 * X2 + u
gen X3_low = 0.05 * X2 + rnormal(0, 2)

quietly regress Y X2 X3_low

* 1. Factor de Inflacion de la Varianza
*    VIF > 10 se considera multicolinealidad severa
estat vif

* 2. Prueba F parcial: H0: alpha_3 = 0
*    Si no se rechaza, X3_low no aporta poder explicativo real
test X3_low

* 3. Criterios de informacion (AIC, BIC): menor valor = modelo preferido
quietly regress Y X2
estat ic

quietly regress Y X2 X3_low
estat ic
\end{lstlisting}

\subsection{Resultados}

\subsubsection{Factor de Inflación de la Varianza}

\begin{lstlisting}[style=stataout]
    Variable |       VIF       1/VIF
-------------+----------------------
          X2 |      1.00    0.999995
      X3_low |      1.00    0.999995
-------------+----------------------
    Mean VIF |      1.00
\end{lstlisting}

Con $X_{3,\text{low}}$ prácticamente ortogonal a $X_2$, el VIF es 1.00,
confirmando que no hay inflación de varianzas en este escenario particular.

\subsubsection{Prueba F parcial sobre $\hat{\alpha}_3$}

\begin{lstlisting}[style=stataout]
 ( 1)  X3_low = 0

       F(  1,   197) =    4.20
            Prob > F =    0.0417
\end{lstlisting}

El estadístico $F = 4.20$ con $p = 0.042$ rechaza $H_0: \alpha_3 = 0$ a
un nivel del 5\%, lo que ilustra que incluso una variable irrelevante puede
parecer significativa por azar en una muestra particular. El resultado es
consistente con $|t| = \sqrt{4.20} \approx 2.05 > 1$, que explica el leve
aumento de $\bar{R}^2$ observado en el escenario de baja correlación.

\subsubsection{Criterios de información AIC y BIC}

\begin{center}
\renewcommand{\arraystretch}{1.3}
\small
\begin{tabular}{lrrrrcc}
\toprule
\textbf{Modelo} & $n$ & $\ell(\text{nulo})$ & $\ell(\text{modelo})$ & $k$ & \textbf{AIC} & \textbf{BIC} \\
\midrule
Correcto: $Y = \beta_1 + \beta_2 X_2 + u$                       & 200 & $-696.47$ & $-287.22$ & 2 & 578.45 & 585.04 \\
Ampliado: $Y = \alpha_1 + \alpha_2 X_2 + \alpha_3 X_{3,\text{low}} + v$ & 200 & $-696.47$ & $-285.11$ & 3 & 576.23 & 586.12 \\
\bottomrule
\end{tabular}
\end{center}

El AIC favorece marginalmente al modelo ampliado ($576.23 < 578.45$), mientras
que el BIC favorece al modelo correcto ($585.04 < 586.12$). El BIC penaliza
más severamente los parámetros adicionales y, en este caso, identifica
correctamente al modelo parsimonioso. Cuando AIC y BIC discrepan, la
literatura recomienda preferir el BIC en muestras medianas a grandes.

%------------------------------------------------------------------------------
\section{Interpretación integrada de los resultados}
%------------------------------------------------------------------------------

Los resultados de Stata confirman con precisión las predicciones teóricas:

\paragraph{Sobre la insesgadez (inciso b).}
En los cuatro modelos estimados, la pendiente oscila entre $2.8618$ y
$2.9858$, convergiendo al verdadero valor $\beta_2 = 3$ con independencia
de la correlación que $X_3$ tenga con $X_2$. El intercepto ronda $2.04$--$2.07$
en todos los casos, próximo al verdadero $\beta_1 = 2$. Los coeficientes de
la variable irrelevante son pequeños: $-0.0809$ (SE $= 0.0395$),
$-0.0110$ (SE $= 0.0496$) y $0.1296$ (SE $= 0.1431$), sin evidencia sólida
de efecto sistemático en los tres primeros escenarios.

\paragraph{Sobre las varianzas (inciso c).}
La inflación de varianzas sigue el patrón del VIF. Con correlación baja
($r_{23} \approx 0.03$), el VIF empírico es $0.98 \approx 1$ y los errores
estándar permanecen prácticamente idénticos ($0.0274$ vs. $0.0276$). Con
correlación moderada ($r_{23} \approx 0.50$), el VIF sube a $2.02$ y el SE
se amplía un 42\% (de $0.0276$ a $0.0392$). Con correlación alta
($r_{23} = 0.977$), el VIF empírico alcanza $23.23$ frente al teórico de
$22.01$ (diferencia de $-0.35$), y el SE se dispara de $0.0276$ a $0.1330$,
un incremento del 382\%.

\paragraph{Sobre $R^2$ y $\bar{R}^2$ (inciso a).}
El $R^2$ convencional crece marginalmente al añadir $X_3$ en todos los
casos (de $0.9833$ a valores entre $0.9833$ y $0.9837$). El $\bar{R}^2$,
en cambio, cae en los escenarios de correlación moderada y alta (de
$0.9832$ a $0.9831$), y sube levemente en el caso de baja correlación
(a $0.9835$) porque el $|t|$ de $X_3$ supera marginalmente la unidad en
esa muestra particular. Esto ilustra que $\bar{R}^2$ no es una garantía
mecánica de caída, sino que depende del estadístico $t$ realizado.

%------------------------------------------------------------------------------
\section{Síntesis: conexión con los ejercicios anteriores}
%------------------------------------------------------------------------------

\begin{center}
\renewcommand{\arraystretch}{1.4}
\small
\begin{tabular}{p{1.0cm}p{4.0cm}p{2.2cm}p{2.6cm}p{3.0cm}}
\toprule
\textbf{Ej.} & \textbf{Error de especificación} & \textbf{Insesgadez} & \textbf{Eficiencia} & \textbf{Validez de inferencia} \\
\midrule
13.2 & Intercepto innecesario incluido      & Preservada  & Reducida            & Válida (menos potente) \\
13.3 & Intercepto relevante omitido         & Perdida     & N/A (hay sesgo)     & Inválida \\
\textbf{13.4} & \textbf{Variable irrelevante $X_3$ incluida} & \textbf{Preservada} & \textbf{Reducida (VIF)} & \textbf{Válida (menos potente)} \\
\bottomrule
\end{tabular}
\end{center}

Los ejercicios 13.2 y 13.4 son cualitativamente análogos (sobre-especificación):
en ambos se preserva la insesgadez pero se pierde eficiencia. La diferencia
cuantitativa es que en 13.2 el coste es un único grado de libertad
desperdiciado (el intercepto), mientras que en 13.4 el daño depende
crucialmente de $r_{23}$: puede ser despreciable si $X_3 \perp X_2$ o
catastrófico si $r_{23} \approx 1$.

%------------------------------------------------------------------------------
\section{Conclusión}
%------------------------------------------------------------------------------

\begin{clave}
  La inclusión de una variable irrelevante en un modelo de regresión
  múltiple tiene consecuencias asimétricas y bien definidas:

  \begin{enumerate}[leftmargin=1.4em]
    \item \textbf{No introduce sesgo}: los estimadores MCO de $\beta_1$ y
      $\beta_2$ mantienen su insesgadez y consistencia. En promedio, el
      modelo sobreajustado recupera los parámetros correctos.

    \item \textbf{Sí produce ineficiencia}: las varianzas se inflan por un
      factor $1/(1 - r_{23}^2)$. Con $r_{23} = 0.977$, como en la
      simulación, la varianza se multiplica por 22 y el error estándar
      crece más de cuatro veces.

    \item \textbf{El $R^2$ convencional es engañoso}: siempre aumenta o
      permanece igual al añadir variables. El $\bar{R}^2$ ajustado
      corrige este sesgo y en general cae cuando $|t_{\hat{\alpha}_3}| < 1$.

    \item \textbf{La inferencia sigue siendo formalmente válida}, pero es
      menos potente: los intervalos de confianza son innecesariamente
      amplios y las pruebas de hipótesis tienen menor capacidad para
      detectar efectos reales.
  \end{enumerate}

  La regla práctica de Gujarati queda completamente justificada: el costo
  de incluir una variable irrelevante es solo ineficiencia, mitigable con
  muestras más grandes; el costo de omitir una variable relevante es sesgo
  sistemático que no desaparece con $n \to \infty$. No obstante, incluir
  variables irrelevantes con alta correlación con los regresores del modelo
  verdadero puede tornar indetectables efectos reales, lo que también es
  problemático en la práctica econométrica.
\end{clave}

\vspace{14pt}
\noindent\rule{\linewidth}{0.4pt}

\noindent{\small\textit{Nota metodológica.}
La asimetría entre sobre-especificación y sub-especificación tiene
implicaciones directas para la estrategia de modelación econométrica.
La \textit{regla general} es partir del modelo más general y reducirlo
solo si la teoría económica y las pruebas estadísticas lo justifican
conjuntamente. Mientras que eliminar variables por parsimonia sin respaldo
teórico introduce el sesgo documentado en el ejercicio~13.3, mantener
variables claramente irrelevantes con alta multicolinealidad puede
producir precisión tan baja que la inferencia carezca de valor práctico.
El equilibrio correcto exige teoría económica sólida como primer criterio,
seguido del diagnóstico estadístico como herramienta auxiliar.}

\vspace{14pt}

\begin{center}
\small\textsc{Referencias}\\[6pt]
\end{center}

\noindent Gujarati, D.~N., \& Porter, D.~C. (2010). \textit{Econometría}
(5.a~ed.). McGraw-Hill. Capítulo~13, secciones~13.2 y~13.3.

\noindent Greene, W.~H. (2018). \textit{Econometric Analysis} (8.a~ed.).
Pearson. Capítulo~5.

\noindent Wooldridge, J.~M. (2016). \textit{Introductory Econometrics}
(6.a~ed.). Cengage. Capítulo~3.

\end{document}
